\documentclass[10pt,a4paper]{amsart}
\usepackage[margin=19mm]{geometry}
\usepackage{amsmath,amssymb,mathtools}
\usepackage[scr=rsfs,cal=euler]{mathalfa}
\usepackage{xfrac}
\usepackage[unicode,hidelinks,bookmarks=false]{hyperref}
\newcommand{\ie}{i.e.\ }
\newcommand{\cf}{cf.\ }
\newcommand{\resp}{respectively\ }
\newcommand{\Subsection}{\subsection}
\providecommand{\nobreakdash}{\nobreak\hbox{-}}
\setlength{\emergencystretch}{2em}
\multlinegap0pt
\newcommand{\gl}{\mathfrak{gl}}
\usepackage{amssymb}
\usepackage{mathtools}
\usepackage{xcolor}
\usepackage{tikz}
\usepackage{dsfont}
\newtheorem{definition}{Definition}
\newtheorem{lemma}{Lemma}
\usepackage{cleveref}
\crefname{definition}{Definition}{Definitions}
\crefname{lemma}{Lemma}{Lemmas}
\newcommand{\EE}{\mathsf{E}}
\newcommand{\F}{\mathrm{F}}
\newcommand{\FF}{{\mathbb{F}}}
\newcommand{\RR}{{\mathbb{R}}}
\newcommand{\hes}{\mathsf{h}}
\renewcommand{\u}{\mathfrak{u}}
\title{Counting the number of elements in the nilradical of a parabolic subalgebra of $\gl_n(\mathbb F_q)$ with a specified Jordan form}
\author{Mohammad Bardestani, Keivan Mallahi-Karai, Samrith Ram, Hadi Salmasian}
\date{Source: arXiv:2503.15440v4 (CC BY 4.0)}
\begin{document}
\maketitle

\noindent\emph{Source and licence.} Counting the number of elements in the nilradical of a parabolic subalgebra of $\gl_n(\mathbb F_q)$ with a specified Jordan form. Version arXiv:2503.15440v4 (CC BY 4.0), distributed by arXiv under the Creative Commons Attribution 4.0 International licence, http://creativecommons.org/licenses/by/4.0/, as recorded in the arXiv licence metadata for this exact version and shown on the abstract page https://arxiv.org/abs/2503.15440v4. Full licence notice: \texttt{licenses/arxiv-2503.15440-CC-BY-4.0.txt}.\par\smallskip

\noindent\textit{Editorial scope.} Counted unit: the article's lemma lem:modularl (source0.tex lines 2323--2326). The article's complete original proof of it is delivered with it (source0.tex lines 2327--2380, one \texttt{proof} environment of 54 lines). No mathematics is written, repaired or re-derived: the statement and the proof are the article's own lines, in the article's own notation, and no retained line is substituted or re-flowed.  Retained uncounted as the source-local context the proof needs: lines 2248--2288; lines 2271--2273; lines 2304--2306; lines 2309--2311; lines 2313--2316.  Nothing was omitted from any retained span, so the package carries no omission record.  No retained line contains a citation, so the wrapper carries no bibliography and no bibliography entry.  No retained line contains a figure environment, an image-inclusion command or drawing code, so no image is delivered and the package declares no asset dependency.  Editorial formatting only, in addition: an AMS \texttt{amsart} preamble that restates the article's own declarations and macros used by the retained lines, namely the environments \texttt{definition}, \texttt{lemma} on separate counters, together with the standard packages those lines need.  The retained environments print in this document as Definition 1, Lemma 1 in order of appearance, whereas the article prints the same items with the numbering of its own full document.  The complete native source of the article as distributed in the arXiv e-print for this exact version is bundled unchanged as \texttt{sources/175059/source0.tex}.
\begin{definition}\label{triple_cond} Let $\hes_0, \hes_1, \hes_2 \in \mathcal{H}_n$ be Hessenberg functions. We say $(\hes_0,\hes_1,\hes_2)$ is a compatible sequence whenever one of the following two conditions holds:

\begin{description}
\item[Case I] There exists $i \in [n-1]$ such that 
\begin{equation}\label{caseI_h}
\hes_1(i-1) < \hes_1(i) < \hes_1(i+1)\quad \text{and}\quad \hes_1(\hes_1(i)) = \hes_1(\hes_1(i)+1).
\end{equation}
Moreover $\hes_0$ and $\hes_2$ are obtained from $\hes_1$ via:
\begin{equation}\label{caseI_modular}
\hes_0(l)= 
\begin{cases}
\hes_1(l) &  l \neq i, \\
\hes_1(i) - 1 &  l = i,
\end{cases} \qquad
\hes_2(l) = 
\begin{cases}
\hes_1(l) & l \neq i, \\
\hes_1(i) + 1 & l = i.
\end{cases}
\end{equation}
\medskip

\item[Case II] There exists $i \in [n-1]$ such that 
\begin{equation}\label{caseII_h}
\hes_1(i+1) = \hes_1(i)+1\quad \text{and}\quad \hes_1^{-1}(i) = \emptyset.
\end{equation}
Moreover $\hes_0$ and $\hes_2$ are obtained from $\hes_1$ via:
\begin{equation}\label{caseII_modular}
\hes_0(l) = 
\begin{cases}
\hes_1(l) &  l \neq i+1, \\
\hes_1(i) &  l = i+1,
\end{cases} \qquad
\hes_2(l) = 
\begin{cases}
\hes_1(l) & l \neq i, \\
\hes_1(i)+1 &  l= i.
\end{cases}
\end{equation}
\end{description}
\end{definition}

\begin{equation}
\hes_1(i+1) = \hes_1(i)+1\quad \text{and}\quad \hes_1^{-1}(i) = \emptyset.
\end{equation}

\begin{equation}\label{f_mu_modular}
    f_\mu: \mathcal{H}_n\to \RR,\qquad \hes\mapsto q^{|h|}\F_{\mu\hes}(q),
\end{equation}

\begin{equation}\label{F_mu_modular}
    (q+1)\F_{\mu\hes_1}(q)=\F_{\mu\hes_0}(q)+q\F_{\mu\hes_2}(q).
\end{equation}

\begin{equation}
\label{elementary_oper}
\mathrm{Row}_l(\EE_{st}(x)A) = \begin{cases} \mathrm{Row}_l(A) & l \neq s; \\ x \,\mathrm{Row}_t(A) + \mathrm{Row}_s(A) & l = s. \end{cases}
\end{equation}

\begin{lemma}
\label{lem:modularl}
 The function $f_\mu: \mathcal{H}_n\to \mathbb{R}$ defined in~\cref{f_mu_modular} satisfies the $q$-modular law. 
\end{lemma}

\begin{proof}
We consider two cases according to~\cref{triple_cond}. 
\medskip

{\noindent \bf Case I:} Set $j=\hes_1(i)$. Define 
$$
\mathcal{A}_0=\{(a_{st})\in \u_{\hes_0}(\FF_q)\cap \mathcal{C_\mu}: a_{ij}\neq 0\}\qquad
\mathcal{A}_1=\{(a_{st})\in \u_{\hes_1}(\FF_q)\cap \mathcal{C_\mu}: a_{i(j+1)}\neq 0\}.
$$
One observes that
$$
\F_{\mu\hes_0}(q)=\F_{\mu\hes_1}(q)+|\mathcal{A}_0|\qquad
\F_{\mu\hes_1}(q)=\F_{\mu\hes_2}(q)+|\mathcal{A}_1|. 
$$
Thus~\cref{F_mu_modular} follows from $q|\mathcal{A}_1|=|\mathcal{A}_0|$. To prove the latter equality, consider 
$\varphi: \mathcal{A}_0\to \mathcal{A}_1$ defined by the assignment
$$
A\mapsto \mathrm{Per}(j,j+1)\,\EE_{j(j+1)}\left(\frac{a_{i(j+1)}}{a_{ij}}\right)\,A\,\EE_{j(j+1)}\left(-\frac{a_{i(j+1)}}{a_{ij}}\right)\,\mathrm{Per}(j,j+1).
$$
We will show that $\varphi(A)\in\mathcal{A}_1$ with $A_{ij}=\varphi(A)_{i(j+1)}$. 
It follows from~\cref{elementary_oper} that the $(i,j+1)$ entry of $A\,\EE_{j(j+1)}\left(-{a_{i(j+1)}}/{a_{ij}}\right)$ is zero. The assumptions $\hes_1(i)<\hes_1(i+1)$ and $\hes_1(\hes_1(i))=\hes_1(\hes_1(i)+1)$ imply that $i+1\leq j$. From~\cref{elementary_oper} we deduce  $A_1=\EE_{j(j+1)}\left({a_{i(j+1)}}/{a_{ij}}\right)\,A\,\EE_{j(j+1)}\left(-{a_{i(j+1)}}/{a_{ij}}\right)$ belongs to $\mathcal{A}_0$ and its $(i,j+1)$ entry is zero. Since $\hes_1(i)<\hes_1(i+1)$ we have $a_{(i+1)(j+1)}=0$. Then $A_1\, \mathrm{Per}(j,j+1)$, which swaps the $j$th and $(j+1)$th column, belongs to $\u_{\hes_1}$ with a non-zero $(i,j+1)$ entry. Thus, the matrix $\mathrm{Per}(j,j+1)\, A_1\, \mathrm{Per}(j,j+1)$ belongs to $\u_{\hes_1}$ with a non-zero $(i,j+1)$ entry, because $i<j$. Since the Jordan form of $A_1$ is $J_\mu$, we conclude that $\varphi(A) = \mathrm{Per}(j,j+1)\, A_1\, \mathrm{Per}(j,j+1) \in \mathcal{A}_1$.


\smallskip

We now show surjectivity of $\varphi$. If $B\in \mathcal{A}_1$ with $b=b_{i(j+1)}\neq 0$, then one can check that for any $x\in\FF_q$ 
\begin{equation}\label{fiber_phi}
\EE_{j(j+1)}(-x/b)\,\mathrm{Per}(j,j+1)\,B\,\mathrm{Per}(j,j+1)\,\EE_{j(j+1)}(x/b)\in \mathcal{A}_0,
\end{equation}
with $(i,j)$ entry $b$ and $(i,j+1)$ entry $x$. This matrix maps to $B$, so $\varphi$ is surjective. Moreover, the matrices of the form~\cref{fiber_phi} are all distinct. This follows from the identity $$\mathrm{Per}(j,j+1)\,\EE_{j(j+1)}(x)\,\mathrm{Per}(j,j+1) = \EE_{(j+1)j}(x)$$
and the fact that $\EE_{(j+1)j}(x) B\, \EE_{(j+1)j}(-x) = B$, where $B \in \mathcal{A}_1$, implies $x = 0$. Thus the fiber of each element has cardinality $q$, and so $q|\mathcal{A}_1|=|\mathcal{A}_0|$.

 
\medskip

{\noindent \bf Case II:} This case is similar to the previous case. Set $j=\hes_1(i)$. Define 
$$
\mathcal{B}_0=\{(b_{st})\in \u_{h_0}\cap\mathcal{C}_\mu: b_{(i+1)(j+1)}\neq 0\},\qquad \mathcal{B}_1=\{(b_{st})\in \u_{h_1}\cap\mathcal{C}_\mu: b_{i(j+1)}\neq 0\}.
$$
Then we have 
$$
\F_{\mu\hes_0}(q)=\F_{\mu\hes_1}(q)+|\mathcal{B}_0|\qquad
\F_{\mu\hes_1}(q)=\F_{\mu\hes_2}(q)+|\mathcal{B}_1|. 
$$
This implies that~\cref{F_mu_modular} follows from $q|\mathcal{B}_1|=|\mathcal{B}_0|$. To prove this we now define
the map 
$\psi: \mathcal{B}_0\to \mathcal{B}_1
$
by the assignment  
$$
B\mapsto \mathrm{Per}(i,i+1)\,\EE_{i(i+1)}\left(\frac{-b_{i(j+1)}}{b_{(i+1)(j+1)}}\right)\,B\,\EE_{i(i+1)}\left(-\frac{b_{i(j+1)}}{b_{(i+1)(j+1)}}\right)\,\mathrm{Per}(i,i+1).
$$
From~\cref{caseII_h} we have $j>i$ which implies $\psi(B)\in \mathcal{B}_1$ and a similar argument as above shows the fiber of each element has cardinality $q$. 
\end{proof}

\end{document}
